\documentclass[11pt,a4paper]{article}
\usepackage[T1]{fontenc}
\usepackage[utf8]{inputenc}
\usepackage[brazil]{babel}
\usepackage[margin=2.2cm]{geometry}
\usepackage{booktabs,graphicx,caption,subcaption,float,amsmath,amssymb,bm,xcolor,array}
\usepackage[output-decimal-marker={,}]{siunitx}
\usepackage[hidelinks]{hyperref}
\graphicspath{{figuras/}}

\title{Refinamento por realização e Karhunen-Loève no talude de Vargas Ceron et al.\ (2025)}
\date{11 de novembro de 2026}
\author{}
\setlength{\emergencystretch}{2em}

% Cada seção fica em secoes/<nome>.tex; sem o arquivo, entra só o título (com o rótulo), e o documento compila.
\newcommand{\secao}[3]{\IfFileExists{secoes/#1.tex}{\input{secoes/#1}}{\section{#2}\label{#3}\textit{(em preparação)}}}

\begin{document}
\maketitle
\tableofcontents

\secao{introducao}{Introdução}{sec:introducao}
\secao{refino_metodo}{O refinamento por realização}{sec:ref:metodo}
\secao{refino_niveis}{Níveis de refinamento e malha inicial mais grossa}{sec:ref:niveis}
\secao{kl_formulacoes}{As formulações de Karhunen-Loève}{sec:kl:form}
\secao{kl_cilindro}{Karhunen-Loève no cilindro de Mika et al.\ (2021)}{sec:kl:cil}
\secao{kl_talude}{Karhunen-Loève no talude}{sec:kl:tal}
\secao{conclusoes}{Conclusões}{sec:conclusoes}

\end{document}
